\section{Related Work}
\label{sec:related}

\subsection{Transit demand forecasting}

Short-term transit demand prediction has been studied extensively. Classical
approaches rely on time series regression or ARIMA on aggregated ridership
data \cite{box1976timeseries}. More recently, graph neural networks have
been used to model spatial dependencies between stops: Li et
al.~\cite{li2023pagtsn} propose PAG-TSN, a parallel attention graph network
for ridership demand that exploits inter-stop adjacency in smart
transportation services. Tree-ensemble baselines (Random Forest, gradient
boosting) remain competitive on tabular transit features, particularly
when combined with classical temporal predictors such as lag terms and
rolling statistics. The ZTBus dataset we use was introduced by Widmer et
al.~\cite{widmer2023ztbus}, which describes the collection protocol and
channel schema but does not present a demand classification pipeline.

\subsection{Time series foundation models}

A recent generation of large pretrained models aims to provide universal
time series forecasting capabilities analogous to language models for text.
TimesFM~\cite{das2024timesfm}, a 200M-parameter decoder-only transformer
from Google Research, and related models such as Chronos, Moirai, and
Lag-Llama follow similar patterns: patch the time series, embed via a
tokenizer, pass through a stack of transformer layers, and project to a
forecast head. Empirical evaluation of these models remains an open
research area: Meyer et al.~\cite{meyer2025benchmarking} survey
benchmarking challenges specific to time series foundation models,
highlighting risks of information leakage through overlapping datasets
and the memorization of global patterns from external shocks. Our work
uses TimesFM not as a forecaster but as a \emph{feature extractor}, a use
case that has received less empirical scrutiny and, as we show, exposes
limitations in the generality of its learned representations when applied
to structured tabular downstream tasks.

\subsection{Feature leakage and evaluation integrity}

Feature leakage, the accidental introduction of target information into
features, is a long-recognized pitfall in applied machine
learning~\cite{kaufman2012leakage}. Kaufman et al.\ formalize leakage as
the inclusion of data in the training process that would not legitimately
be available at prediction time, enumerate common sources (time-based
leakage, target encoding, group-based splits), and recommend a set of
detection and avoidance strategies. In time series contexts specifically,
leakage can arise from train-test temporal overlap, look-ahead bias in
rolling statistics, or cross-series correlation at the benchmark
level~\cite{meyer2025benchmarking}. Our contribution identifies a
different class of leakage: not at evaluation time, but at \emph{feature
engineering} time, introduced by an interaction between imputation and
feature computation that appears benign when the two steps are examined
separately.

\subsection{Sparse event time series}

Sparse event data is common in many domains: medical vital signs recorded
only during checkups, IoT sensor readings that trigger on thresholds, and
financial tick data between quotes. Bagattini et
al.~\cite{bagattini2019sparse} propose a classification framework
specifically for sparse multivariate temporal features in medical records,
introducing a three-phase symbolic transformation that handles variable
sparsity without imputation. Dillmann et al.~\cite{dillmann2024events}
develop sparse autoencoder representations of event time series for
high-energy astrophysics, operating directly on energy-time binned event
files without intermediate regularization. Che et
al.~\cite{che2018grud} introduce GRU-D, which directly models
missing-value patterns in multivariate healthcare time series rather than
imputing at all. All of these approaches address sparsity by building
models that treat sparsity as a first-class signal. We take the opposite
approach: we use forward-fill imputation for its simplicity and physical
grounding, then fix the resulting leakage at the feature engineering
stage rather than the model.

\subsection{Classical time series analysis}

Our work makes extensive use of classical time series tools. Box and
Jenkins~\cite{box1976timeseries} provide the foundational ARIMA framework
we use for our temporal-only baseline. Dickey and
Fuller~\cite{dickey1979unit} and Kwiatkowski et
al.~\cite{kwiatkowski1992kpss} provide the ADF and KPSS unit root tests,
which we apply in complementary fashion to establish stationarity
properties. Autocorrelation (ACF) and partial autocorrelation (PACF)
diagnostics remain the standard tools for characterizing temporal memory
in a series.
